\documentclass[a4paper,12pt,reqno]{amsart}

\usepackage{amsmath, amsfonts, amssymb, amsthm, mathrsfs, mathtools, mathalfa, setspace}
\usepackage{caption}
\usepackage{comment}
\usepackage{yfonts}
\usepackage{enumerate}
\usepackage{enumitem}
\usepackage{stmaryrd}
\usepackage{fancyhdr}
\usepackage{bm}
\usepackage[utf8]{inputenc}
\usepackage[pdfencoding=auto,hyperfootnotes=false,backref=page]{hyperref}
\usepackage{tikz,tikz-cd}
\usepackage[hang,flushmargin]{footmisc}
\usepackage{xr}
\usepackage{extarrows}
\usepackage{todonotes}
\usepackage{faktor}
\usepackage{xfrac}
\usepackage{times}
\usepackage{color}
\usepackage{graphicx}
\IfFileExists{quiver.sty}{\usepackage{quiver}}{}
\usetikzlibrary{babel}
\usepackage{xurl}
\usepackage{bbm}


\usetikzlibrary{arrows}
\usetikzlibrary{positioning, arrows.meta}

\makeatletter
\def\@tocline#1#2#3#4#5#6#7{\relax
  \ifnum #1>\c@tocdepth % then omit
  \else
    \par \addpenalty\@secpenalty\addvspace{#2}%
    \begingroup \hyphenpenalty\@M
    \@ifempty{#4}{%
      \@tempdima\csname r@tocindent\number#1\endcsname\relax
    }{%
      \@tempdima#4\relax
    }%
    \parindent\z@ \leftskip#3\relax \advance\leftskip\@tempdima\relax
    \rightskip\@pnumwidth plus4em \parfillskip-\@pnumwidth
    #5\leavevmode\hskip-\@tempdima
      \ifcase #1
       \or\or \hskip 1em \or \hskip 2em \else \hskip 3em \fi%
      #6\nobreak\relax
    \dotfill\hbox to\@pnumwidth{\@tocpagenum{#7}}\par
    \nobreak
    \endgroup
  \fi}
\makeatother

\special{!
/setdistillerparams where {pop}{userdict /setdistillerparams {pop} put}ifelse
 <</AutoRotatePages /None>> setdistillerparams
}


\special{!
/setdistillerparams where {pop}{userdict /setdistillerparams {pop} put}ifelse
 <</AutoRotatePages /None>> setdistillerparams
}

\topmargin -0.2 cm
%\linespread{1.3}
\evensidemargin 0cm
\oddsidemargin 0cm
\textheight 24 cm
%\textheight 615pt
\marginparwidth = 55pt
\textwidth 16cm
%\textwidth 360pt
\setlength{\parindent}{0.7 cm}
\setlength{\footskip}{1.5 cm}
\setlength {\marginparwidth }{2cm}

\renewcommand{\baselinestretch}{1.5}
\newcommand{\myp}{{\vskip0.2truein}}
\newcommand{\fixme}[1]{\textcolor{red}{\textbf{(#1)}}}
\newcommand{\cubes}{\mathbf{Cubes}}
\newcommand{\Hom}{\mathrm{Hom}}

\newcommand*{\compl}[1]{#1^\dagger}

% ---------- Basic notation ----------
\newcommand{\Gammaact}{T} % action T^\gamma
\newcommand{\cube}[1]{\{0,1\}^{#1}}
\newcommand{\abs}[1]{\left|#1\right|}
\newcommand{\1}{\mathbf{1}}
\newcommand{\0}{\mathbf{0}}
\newcommand{\del}{\partial}
\newcommand{\cH}{\mathcal{H}}


% parity / alternating signs on cube vertices
\newcommand{\sgn}[1]{(-1)^{|#1|}}


\newtheorem{theorem}{Theorem}[section]
\newtheorem{lemma}[theorem]{Lemma}
\newtheorem{definition}[theorem]{Definition}
\newtheorem{corollary}[theorem]{Corollary}
\newtheorem{proposition}[theorem]{Proposition}
\newtheorem{question}[theorem]{Question}
\newtheorem{problem}[theorem]{Problem}
\newtheorem{conjecture}[theorem]{Conjecture}
\theoremstyle{definition}
\newtheorem{defn}[theorem]{Definition}
\newtheorem{remark}[theorem]{Remark}
\newtheorem{example}[theorem]{Example}

\newtheorem{assumption}[theorem]{Assumption}
\newtheorem{goal}[theorem]{Goal}


\newcommand{\R}[0]{\hbox{$\mathcal R$}}
\newcommand*{\sbr}[1]{\scalebox{0.8}{$(#1)$}}
\newcommand*{\db}[1]{\llbracket #1\rrbracket}
\newcommand{\tb}{\textbf}
\newcommand{\mk}{\mathfrak}
\newcommand{\mc}{\mathcal}
\newcommand{\ms}{\mathscr}
\newcommand{\mt}{\mathtt}
\newcommand{\mf}{\mathbf}
\newcommand{\mb}{\mathbb}
\newcommand{\mr}{\mathrm}
\newcommand{\parc}{\partial}
\newcommand{\wh}{\widehat}
\newcommand{\wt}{\widetilde}
\newcommand{\st}{\mathrm{st}}
\newcommand{\nst}[1]{{^*#1}} % non-standard
\renewcommand{\theenumi}{\roman{enumi}}
\renewcommand{\labelenumi}{\textup{(}\theenumi\textup{)}}
\newcommand{\eulerian}[2]{\genfrac{\langle}{\rangle}{0pt}{}{#1}{#2}}
\newcommand{\ud}{\,\mathrm{d}}
\newcommand{\id}{\mathrm{id}}


\newcommand{\Var}{\mathrm{Var}}
\newcommand{\stab}{\mathrm{Stab}}
\newcommand{\aut}{\mathrm{Aut}}

\newcommand{\ab}{\mathrm{Z}} 



%%%%%%%%%%%%%%%%%%%%%%%%%%%
\newcommand{\Bo}{\mathrm{B}}
\newcommand{\tran}{\mathrm{\Theta}}
\newcommand{\hcf}{\mathrm{hcf}}
\newcommand{\lcm}{\mathrm{lcm}}
\newcommand{\ia}{\wp}
\newcommand{\poly}{\mathrm{poly}}
\newcommand{\Span}{\mathrm{Span}}
\newcommand{\Supp}{\mathrm{Supp}}
\newcommand{\Lip}{\mathrm{Lip}}
\newcommand{\adj}{\mathrm{adj}}
\newcommand{\Id}{\mathrm{Id}}
\newcommand{\sign}{\mathrm{sign}}
\newcommand{\codim}{\mathrm{codim}}
\newcommand{\comp}{\mathrm{K}}
\newcommand{\coup}{\mathsf{Cg}}
\newcommand{\q}{\mathrm{c}}
\newcommand{\ns}{\mathrm{X}}
\newcommand{\nss}{\mathrm{Y}}
\newcommand{\co}{\circ\hspace{-0.02cm}}
\newcommand{\cu}{\mathrm{C}}
\newcommand{\cor}{\mathrm{Cor}}
\newcommand{\cs}{\mathrm{s}}
\newcommand{\upmod}{\perp\!\!\!\perp}
\newcommand{\Aff}{\mathrm{Aff}}
\newcommand{\bnd}{\mathrm{B}}
\newcommand{\abph}{\mathrm{U}}
\newcommand{\diam}{\mathrm{diam}}
\newcommand{\nor}{\mathrm{\Xi}}
\newcommand{\mcd}{\mathrm{mcd}}
\newcommand{\inv}{\mathrm{inv}}
\newcommand{\orb}{\mathrm{orb}}
\newcommand{\ev}{\mathrm{ev}}
\newcommand{\tRe}{\mathrm{Re}}
\newcommand{\tIm}{\mathrm{Im}}
\newcommand{\Zmod}[1]{\Z/#1\Z} % \Z/#1\Z or \Z_{#1} -- easily change from Z_p to Z/pZ

%%%%%%%%%%%%%%%%%%%%%%%%%%%%%

\providecommand{\abs}[1]{\lvert#1\rvert}
\providecommand{\arr}[1]{\langle#1\rangle}
\providecommand{\norm}[1]{\lVert#1\rVert}
\providecommand{\Tnorm}[1]{\lVert#1\rVert_{\T}}
\providecommand{\md}{d}
\providecommand{\GL}[1]{\GLs(#1,2)\rtimes \mathbb{Z}_2^{#1}}
\providecommand{\comm}[1]{[#1]}

\setlength {\marginparwidth }{2cm}


\setcounter{tocdepth}{2}


\DeclareMathOperator{\tnss}{\tilde{\nss}}

\begin{document}
\title{On $m$-sum-free sets in $\mb{Z}/p\mb{Z}$}

\author{}

\maketitle

\section{Introduction}

A motivation for this project is to study whether methods involving Linear Programming, used in recent works such as \cite{LMPV}, can help to improve the best known upper bounds for the densities of so-called \emph{$m$-sum-free sets} in cyclic groups of prime order. 

For a positive integer $m$, a set $A\subset \mb{Z}/p\mb{Z}$ is said to be \emph{$m$-sum-free} if there is no triple $(x,y,z)\in A^3$ satisfying the equation $x+y=mz$. Equivalently, the set $A$ is $m$-sum-free if the sum of dilates $A+A-m\cdot A$ does not contain $0$. Also equivalently, $A$ is $m$-sum-free if $(A+A)\cap (m\cdot A) =\emptyset$.

Using the Cauchy-Davenport theorem and this last equivalent definition, it is readily seen that the density of an $m$-sum-free set is at most 1/3. Interestingly, this natural upper bound is far from sharp, and it is a non-trivial problem to improve the upper bounds on the density of $m$-sum-free sets (below we give more details on known results in this direction).

The topic of $m$-sum-free sets can be connected with linear programming methods such as those of \cite{LMPV} as follows:  Theorem 1 from \cite{LMPV} immediately implies that a ``$2$-sum-free" set in $\mb{Z}/p\mb{Z}$ must have density at most $2/7$ (which is less/better than the ``trivial" Cauchy-Davenport bound of $1/3$ mentioned above). Of course, in this $m=2$ case of the $m$-sum-free-set problem, this upper bound is not interesting.\footnote{Strictly speaking, a non-empty set $A$ cannot be 2-sum-free since any $x\in A$ yields a solution $(x,x,x)$ to the equation that $A$ is supposed to avoid; even if we strengthened the $2$-sum-free property by requiring a solution $(x,y,z)$ to have \emph{distinct} coordinates $x,y,z$, still Roth's theorem would tell us that $|A|$ must in fact be at most $o(p)$.} However, if the linear programming methods used in \cite{LMPV} could be extended to generalize \cite[Theorem 1]{LMPV} for $m>2$, then this would give upper bounds for the densities of sets $A$ such that  $A+A-m\cdot A$ does not fill all of $\mb{Z}/p\mb{Z}$. This could yield new upper bounds which, for $m>2$, would be interesting.

\section{The state of the art on bounds for $m$-sum-free sets}

The best bounds on $m$-sum-free sets for $m\geq 3$ that are explicitly worked out in the literature are (to the best of our knowledge) in \cite{CGG}, namely Theorems 1.5 and 1.7 in that paper, which together tell us that the maximum density of an $m$-sum-free set is between $1/8$ and $1/3.1955$. 

Note that the $2/7$ upper bound implied by \cite[Theorem 1]{LMPV} is less than $1/3.1955$, so a generalization of \cite[Theorem 1]{LMPV} to $m\geq 3$ with a similar bound close to 2/7 would already give an improvement on this result from \cite{CGG}. 

However, there is more recent work \cite{Gryn} by Grynkiewicz which improves a central tool that was used in \cite{CGG} to get the upper bound $1/3.1955$ (that tool being \cite[Theorem 1.4]{CGG}, i.e.\ a Freiman $3k-4$ type result in $\mb{Z}/p\mb{Z}$ optimized for mid-density sets). At first sight, Grynkiewicz's new main result (i.e. Theorem 1.3 in \cite{Gryn}) does not improve directly upon \cite[Theorem 1.4]{CGG}. However, the paper \cite{Gryn} is very technical and full of variants of its main result, and it is not clear to us yet that some such variant (e.g. \cite[Theorem 3.3]{Gryn}) cannot be optimized easily to improve the upper bound 1/3.1955 to some smaller number. We would need to check/calculate this to know what the LP method would really have to beat. This is work in progress which we will try to report upon in this section soon.

\bigskip

\paragraph{Update (three-parameter LP).}
After the writeup of Section \ref{sec:initial-obstacles} below we found that
the obstruction can actually be overcome by letting the witness forms depend
on three parameters $(u,v,w)$ rather than one. With that generalization the
density bound for $m=3$ drops to $16/33\approx 0.4848$ with only six witness
forms, to $48/103\approx 0.4660$ with nine forms, and to approximately
$0.4398$ with twelve forms. The same construction works uniformly for every
$m\ge 3$. A full account is given in \texttt{latex/6\_three\_parameter\_model.tex}
and the numerical table is in \texttt{latex/7\_three\_parameter\_results.tex}.
These bounds are still above the unconditional $1/3.1955$ bound of
\cite{CGG}, so the LP approach is not yet competitive with the
additive-combinatorial machinery, but they are the first LP relaxation bounds
that strictly break the $1/2$ barrier for $m\ge 3$.

\bigskip

\todo[inline]{Please feel free to add anything here, e.g. new sections to explain or explore the possible extensions of the LP method. (Pablo)}

\section{Overview of the LP method}

The approach in \cite{LMPV} is to view the density of a set $A \subset \mb{Z}/p\mb{Z}$
as a linear function of the densities of related \emph{atomic} sets $S_{\epsilon} = \bigcap_{i \in I}(A - x_i)^{\epsilon_i}$, for $\epsilon \in \{-1, 1\}^{\abs{I}}$ and taking $X^{-1}$ to be $\mb{Z}/p\mb{Z} \setminus X$. These \emph{atomic densities} must satisfy linear relations such as being nonnegative, having a sum of $1$, and translation invariance in the sense that if $j_k = i_k + t$ for some choices of $i_k, j_k \in [1, \abs{I}]$, then $\delta\left(\cap_k(A - x_{i_k})\right) = \delta\left(\cap_k(A - x_{j_k})\right)$.

Furthermore, assume that $d = 6 \notin A + A - 2 \cdot A$. If $\epsilon_i = \epsilon_j = \epsilon_k = 1$ and $x_i + x_j - 2x_k = d$, then $S_\epsilon$ must be empty, making its density 0. If we view the densities of all the sets as variables in a linear program, we can maximize the density of $A$. A certain small choice of $x_i$ that does not depend on $p$ yields the upper bound $2/7$.

\section{Initial obstacles for generalizing the LP method}\label{sec:initial-obstacles}

One might think that this argument easily translates to our setup since we are trying to find the maximum density of a set $A$ such that $0 \notin A + A - k \cdot A$. However, trying to apply it directly, one runs into the following issue. Suppose $c \in (A - x_1) \cap (A - x_2) \cap (A - x_3)$. That is, $x_1 = a_1 - c$, $x_2 = a_2 - c$, and $x_3 = a_3 - c$ for $a_1, a_2, a_3 \in A$. Since $1 + 1 - k = 2 - k \neq 0$, there is no way to combine these expressions into $a_1 + a_2 - ka_3 = 0$ with the $c$'s canceling out.

One possible fix to this is to explore other choices of the atomic sets. For example, for $k = 4$, one might consider the sets $S_{\epsilon}$ as before, $T_{\mu} = \bigcap_{j \in J}(2 \cdot A - y_j)^{\mu_j}$, and the atoms $U_{\epsilon, \mu} = S_{\epsilon}\cap T_{\mu}$. Indeed,
if $c \in (A - x_1) \cap (A - x_2) \cap (2 \cdot A - y_3)$ and $x_1 + x_2 - 2y_3 = 0$ then
\[
    0 = c + c - 2c = (a_1 - x_1) + (a_2 - x_2) - 2(2a_3 - y_3) = a_1 + a_2 - 4a_3 - (x_1 + x_2 - 2y_3).
\]
Since $x_1 + x_2 - 2y_3 = 0$, this yields $a_1 + a_2 - 4a_3 = 0$, which contradicts our assumption about the set $A$. This adds similar constraints to our linear program. However, with this approach and after some computer search (\url{https://github.com/ferranEspuna/lp-density-estimates}), I [Ferran] have not been successful in establishing any upper bounds below the trivial $1/2$.
The reason for this is that there is a corresponding trivial solution which we can describe as follows. Let $\epsilon_1 = \mathbbm{1}_{I}$ and $\mu_{1} = \mathbbm{1}_{J}$. Then, setting $\delta(U_{\epsilon_1, -\mu_1}) = \delta(U_{-\epsilon_1, \mu_1}) = 1/2$ and the rest of atomic densities to $0$ satisfies all constraints that we can naively impose with shifts that do not depend on $p$.



\begin{thebibliography}{1}

\bibitem{LMPV} V. Lev, M. Matolcsi, P. P. Pach, D. Varga,   \emph{On the density of kravitz sets}, \url{https://arxiv.org/pdf/2510.16522}.

\bibitem{CGG} P. Candela, D. González-Sánchez, D. Grynkiewicz, \emph{On sets with small sumset and $m$-sum-free sets in $\mb{Z}/p\mb{Z}$}, 
\url{https://arxiv.org/pdf/1909.07967}

\bibitem{Gryn} D. Grynkiewicz, \emph{The $3k-4$ theorem modulo a prime: high density for $A+B$}, 
\url{https://arxiv.org/pdf/2402.15028}

\end{thebibliography}
\end{document}
